% =====================================================================================================================
% Introduction (owner: intro; integrated 2026-10-04). Pre-results draft.
%  * Every result is a \DATANEEDED slot. The only empirical facts stated are the registered gate outcomes of the pilot that
%    motivated the registered remedy (ML-1M zero-shot UAUC 0.587 < 0.60; sports next-item NDCG@10 0.2093 >= 0.1863;
%    KuaiRec zero-shot UAUC 0.469 < 0.58). Nothing else from the pilot is stated.
%  * Branches (remedy outcome, fine-tuning gate outcome, nested method slot) are separate \DATANEEDED{branch: ...} sentences:
%    keep the one that occurred and delete the other.
%  * Findings written into slots must satisfy the registered claim-admission rule (confirmed on untouched users, same sign on at
%    least one fresh Amazon domain and on the second backbone); otherwise the sentence stays direction-free.
%  * Design statements (Qwen3-8B zero-shot and LoRA on four rated panels, Llama-3.1-8B-Instruct as replication backbone,
%    zero-shot only on the four next-item panels, the three actionability routes, five-seed pruning) describe the registered
%    grid (Amendments 2 and 3): re-verify them against what was actually run (cut rules) before submission.
%  * Wording rule (Amendment 3, section 5): do not attribute the item prior to name recognition unless the pseudonym knockout is
%    POSITIVE in both domains and both regimes; until then write "popularity-linked".
%  * Integration (2026-10-04): "registered before any confirmatory result was analysed" is exact (the next-item scoring started
%    before its analysis specification was written; the specification was frozen before any audit score was analysed; Section
%    sec:protocol says so). Holm families and the two single tests follow Amendment 3 section 11. The reported decomposition
%    statistic is the non-prior share, an upper bound on the personal share. The RQ list was folded into one paragraph and the
%    gap paragraph shortened for the page budget (Table tab:rq carries the endpoints).
%  * Cross-references used: sec:prelim, sec:protocol; tab:rq (audit protocol); sec:exp-gates, tab:guide (findings section).
% =====================================================================================================================
\providecommand{\method}{\textsc{Mirror}}
\providecommand{\py}{P(\texttt{Yes})}
\providecommand{\introlistsetup}{\setlength{\leftmargin}{2.1em}\setlength{\labelwidth}{1.9em}\setlength{\labelsep}{0.2em}%
  \setlength{\itemsep}{0pt}\setlength{\parsep}{0pt}\setlength{\topsep}{2pt}\setlength{\partopsep}{0pt}%
  \def\makelabel##1{\textbf{##1}\hfil}}

\section{Introduction}
\label{sec:intro}

Large language models (LLMs) are increasingly used as \emph{pointwise} recommenders: asked ``Will this user like this item?
Answer Yes or No,'' the model's probability of \emph{Yes} is the ranking score~\citep{bao2023tallrec,zhang2025collm,zhang2024binllm}.
The same number is read as a \emph{confidence}: a system may abstain or truncate a list when it is
low~\citep{fan2026ugr,kweon2024perk}, what is shown feeds back into what is learnt, a loop known to amplify popularity
bias~\citep{chaney2018feedback,mansoury2020feedback}, and fine-tuning data can be pruned by uncertainty~\citep{wu2025unit} or
coreset scores~\citep{lin2024dealrec,mei2025goracs}. Whether the confidence deserves this reading is largely untested for the
yes/no family, evaluated by discrimination (AUC, or per-user AUC, UAUC) and, to our knowledge, not by calibration; uncertainty
work on LLM recommenders addresses list-wise rankings~\citep{kweon2025uqrec,yan2026evirank}, generative
recommenders~\citep{fan2026ugr} or verbalised confidence~\citep{ravikumar2026hallucinating}. Missing are labels that match the
question, disaggregation (a ranking depends on the order within one user's candidates, a serving decision on the order across
users) and non-LLM reference points on the same pairs~\citep{ferraridacrema2019progress}.

Above all, it is open how much of $\py$ is a prior about the item and how much is evidence about this user. Write the log-odds
of \emph{Yes} as $\ell(u,i)=b_u+\pi(i)+e(u,i)$: a user offset, an item prior that the model brings to the item whoever asks
(possibly popularity-linked), and personal evidence. Only $e$ personalises the ranking, and the prior is what per-user
calibration cannot touch: by an invariance lemma (Section~\ref{sec:prelim}) any strictly increasing per-user map, such as
temperature or Platt scaling~\citep{guo2017calibration}, leaves per-user AUC and within-user NDCG unchanged; outcomes change only
through item-dependent corrections, cross-user decisions and training.

We therefore audit the confidence of yes/no recommenders through six practitioner questions, S1--S6, posed as research questions
with registered endpoints (Table~\ref{tab:rq}): is the model sure when it is right (RQ1) and unsure when it is wrong (RQ2); does
serving the most confident items over-expose head items (RQ3, static); how well does confidence track labels that match the
question, and how much of it is item prior (RQ4); is it higher for popular items, and is that warranted once quality is
controlled (RQ5); and can uncertainty flag noisy training examples better than matched random pruning (RQ6)?

\paragraph{A registered audit.}
Panels, endpoints, thresholds and a claim-admission rule were registered before any confirmatory result was analysed
(Section~\ref{sec:protocol}): a finding enters the abstract only if it is confirmed on untouched users and has the same sign on
at least one fresh Amazon domain and on a second backbone; until then its wording stays direction-free. Not every gate was met:
under the registered zero-shot prompt a pilot on 1,500 ML-1M users reached UAUC 0.587 against a bar of 0.60 while the next-item
component was met (NDCG@10 0.2093 against 0.1863), and on the fully observed KuaiRec matrix~\citep{gao2022kuairec} the zero-shot
UAUC was 0.469 against 0.58, so exposure is analysed statically. These registered negative results (Section~\ref{sec:exp-gates})
motivated one registered prompt-remedy round. It selected, on development users, the label-aligned prompt that asks whether the
user would rate the item 4 stars or higher and, on 1,683 untouched ML-1M users, that prompt reached UAUC
0.603 (95\% CI 0.594--0.612) against 0.576 for the registered prompt on the same users, so the bar was
met.
As TALLRec-style recommenders are fine-tuned~\citep{bao2023tallrec}, the LoRA regime was to carry the main results if it met its own
registered gate, a mean UAUC of at least 0.65 over three seeds. It did (mean 0.742;
seeds 0.740, 0.745, 0.743), so the fine-tuned regime is the main object and the zero-shot regime its contrast.

\paragraph{Contributions.}
\begin{list}{}{\introlistsetup}
\item[C1] \emph{Protocol.} A pre-registered audit of Qwen3-8B, zero-shot and with LoRA (three seeds), replicated with
  Llama-3.1-8B-Instruct, on rated panels with matched labels (ML-1M; Amazon Toys, Video Games, Sports) and, zero-shot, on
  next-item panels of an existing benchmark protocol (four domains $\times$ 10k users $\times$ 101 candidates), with non-LLM
  references, user-clustered intervals, Holm correction within registered endpoint families (the primary fine-tuning hypothesis
  and the pruning contrast are single tests) and fresh-user confirmation.
\item[C2] \emph{Decomposition.} A reliability-corrected non-prior share (an upper bound on the personal share), the information
  gain over the prior-only item mean, a star-permutation contrast, a two-axis popularity analysis~\citep{xiong2023procal}
  partialled on item quality and, in the fine-tuned programme, a pseudonym knockout with a real-brand placebo.
\item[C3] \emph{Actionability.} The lemma and a controlled test of each route it leaves open: item-dependent corrections (including
  a two-view contrast~\citep{burns2023ccs} that we evaluate, not propose) against a paraphrase placebo and an ensemble null;
  selective serving~\citep{geifman2017selective} against random coverage; exposure against the pool and eight published
  recommenders; pruning against matched random pruning over five seeds.
\item[C4] \emph{A guide.} A one-table answer to S1--S6 (Table~\ref{tab:guide}), each tied to its registered endpoint.
\end{list}
% METHOD-SLOT-BEGIN (introduction): delete this sentence if the nested method slot is retired at its kill date
\DATANEEDED{branch: only if the nested method slot (prior-offset LoRA) survives its kill rule -- one sentence with its gain over
post-hoc stacking, CI and datasets; else omit}
% METHOD-SLOT-END (introduction)

\paragraph{Findings in brief.}
\DATANEEDED{RQ1, RQ2, RQ4: zero-shot versus LoRA confidence--label relation (UAUC, calibration, error-detection AUROC) beside the
reference rows; both backbones; confirmation users; 95\% CI; source: main reliability table; direction-free unless admitted}
\DATANEEDED{RQ4, RQ5: non-prior share (an upper bound on the personal share), information gain over the prior-only item mean,
quality-partialled popularity link, knockout contrast against the placebo; source: decomposition and knockout tables}
\DATANEEDED{RQ3, RQ6, C3: corrections against placebo and ensemble null; selective serving against random; head-share excess over
the pool and against the published recommenders; pruning against matched random; source: actionability tables}
